\documentclass[11pt]{amsart}

\usepackage[T1]{fontenc}
\usepackage{lmodern}
\usepackage{microtype}
\usepackage{mathtools,amssymb,amsthm}
\usepackage[hidelinks]{hyperref}

\hypersetup{
  pdftitle={A nine-vertex counterexample to the weighted Graham pebbling
            conjectures as stated in the e-print of Herscovici, Hester and
            Hurlbert}
}

\newtheorem{theorem}{Theorem}
\newtheorem{lemma}[theorem]{Lemma}
\newtheorem{corollary}[theorem]{Corollary}
\theoremstyle{definition}
\newtheorem{conjecture}{Conjecture}

\DeclareMathOperator{\VAL}{VAL}

\title[A weighted triangle squared]
{A nine-vertex counterexample to the weighted Graham pebbling conjectures
as stated in the e-print of Herscovici, Hester and Hurlbert}

\date{}

\subjclass[2020]{05C57, 05C76, 05C22}
\keywords{graph pebbling, weighted graph, Cartesian product, Graham's
conjecture, pebbling number, counterexample}

\begin{document}

\begin{abstract}
In a weighted graph every edge $e$ carries a weight $w(e)\in\mathbb Z_{>0}$,
and a pebbling move along $e=(x,x')$ removes $w(e)$ pebbles from $x$ and
places one of them on $x'$. The e-print of Herscovici, Hester and Hurlbert
conjectures that $\pi(G\times H,(x,y))\le\pi(G,x)\pi(H,y)$ for all weighted
graphs $G,H$ and all vertices $x\in V(G)$, $y\in V(H)$. Let $G$ be the
triangle on
$\{1,2,3\}$ with $w(1,2)=2$, $w(1,3)=3$, $w(2,3)=2$. Then $\pi(G,v)=3$ at
every vertex $v$, while in $G\times G$ the distribution with seven pebbles on
$(3,3)$ and one pebble each on $(2,3)$ and $(3,2)$ cannot place a pebble on
$(1,1)$; hence
\[
  \pi(G\times G,(1,1))\ \ge\ 10\ >\ 9\ =\ \pi(G,1)\pi(G,1).
\]
The same nine-vertex object refutes the weighted
$st$-pebbling conjecture of that e-print, by its own specialisation $s=t=1$,
and, because $\pi(G)=3$, also the unrooted form $\pi(G\times
H)\le\pi(G)\pi(H)$, which Sieben had asked about. Every statement refuted
here is quoted from the e-print source; the printed text in \emph{Discrete
Math.} was not accessible to us, so it is the e-print formulations that are
refuted. The unsolvability is proved by hand in two cases; a finite
computation adds the exact value $\pi(G\times G,(1,1))=10$. The weights are
non-uniform, so Graham's conjecture is untouched.
\end{abstract}

\maketitle

\section{The conjectures}

We use the definitions of the e-print source of Herscovici, Hester and
Hurlbert \cite{HHH} verbatim. A \emph{weighted graph} is a graph $G=(V,E)$ together
with a function $w:E(G)\to\mathbb N^{+}$, and a \emph{pebbling move along the
edge $e=(x,x')$} consists of removing $w(e)$ pebbles from $x$, moving one of
the pebbles onto $x'$, and throwing the other pebbles away.
A \emph{distribution} $D$
assigns a nonnegative integer $D(v)$ to each vertex, $|D|=\sum_vD(v)$, and
$D$ is \emph{$r$-solvable} if some sequence of moves reaches a distribution
with at least one pebble on $r$. As in the unweighted theory, $\pi(G,r)$ is
the least $N$ such that every distribution of $N$ pebbles is $r$-solvable,
$\pi_t(G,r)$ is the least $N$ such that $t$ pebbles can be moved to $r$ from
every distribution of $N$ pebbles, and
\begin{equation}\label{eq:unrooted}
  \pi(G)=\max_{v\in V(G)}\pi(G,v).
\end{equation}
The \emph{weighted Cartesian product} $G\times
H$ has the vertex set and the edge set of the unweighted product, with
\[
  w\bigl((x,y),(x,y')\bigr)=w(y,y'),\qquad
  w\bigl((x,y),(x',y)\bigr)=w(x,x').
\]
Setting every weight to $2$ recovers
unweighted pebbling, so each weighted conjecture is stronger than its
unweighted counterpart. Finally, any connected weighted graph may be regarded
as a complete graph in which every weight is the least product of weights
along a path between its endpoints; we call such a graph \emph{tight}, and the
graph we exhibit is tight. For the wider theory of pebbling numbers see
Hurlbert's survey \cite{Hurlbert}.

The two statements of the e-print source of \cite{HHH} that concern us are
the following, quoted word for word.

\begin{conjecture}[\cite{HHH}]\label{conj:W}
For all weighted graphs $G$ and $H$ and all vertices $x\in V(G)$ and $y\in
V(H)$, we have $\pi(G\times H,(x,y))\le\pi(G,x)\pi(H,y)$.
\end{conjecture}

\noindent
(E-print source, lines 690--694.)

\begin{conjecture}[\cite{HHH}]\label{conj:Wst}
For all weighted graphs $G$ and $H$, all positive integers $s$ and $t$, and
all vertices $x\in V(G)$ and $y\in V(H)$, we have
\[
  \pi_{st}(G\times H,(x,y))\le\pi_s(G,x)\pi_t(H,y).
\]
\end{conjecture}

\noindent
(E-print source, lines 683--688.)

Conjecture~\ref{conj:Wst} implies Conjecture~\ref{conj:W} by letting
$s=t=1$, as the e-print itself notes (its lines 741--743); this is all we use
of it. The unrooted weighted statement
\begin{equation}\label{eq:WG}
  \pi(G\times H)\ \le\ \pi(G)\pi(H)
  \qquad\text{for all weighted graphs }G,H,
\end{equation}
which is implied by Conjecture~\ref{conj:W} through \eqref{eq:unrooted}, is
among the statements \cite{HHH} says can be made for weighted graphs, and it
is exactly what Sieben asks in the last section of \cite{Sieben}: ``What
general results are there about the pebbling number of weighted graphs? In
particular, does Graham's conjecture hold for weighted graphs?''
(\cite[e-print source, lines 1016--1021]{Sieben}).

On the positive side, Chung essentially proved
Conjecture~\ref{conj:W} when $G$ and $H$ are powers of $K_2$, that is,
cubes in which the weights of parallel edges are equal
(\cite[lines 696--699]{HHH}, citing \cite[Theorem~3]{Chung}).

\section{The counterexample}

Let $G$ be the complete graph on $\{1,2,3\}$ with
\[
  w(1,2)=2,\qquad w(1,3)=3,\qquad w(2,3)=2 .
\]
It is tight, since $2\le3\cdot2$, $3\le2\cdot2$ and $2\le2\cdot3$. For a root
$r$ let $d_r(v)$ denote the least product of weights along a path from $v$ to
$r$, with $d_r(r)=1$; thus $d_1=(1,2,3)$, $d_2=(2,1,2)$ and $d_3=(3,2,1)$.

Two elementary facts are used throughout. First, solvability is monotone:
if $D\le D'$ pointwise and $D$ is $r$-solvable then so is $D'$. Second, if
$D(v)\ge d_r(v)$ for some $v$ then $D$ is $r$-solvable, because $d_r(v)$
pebbles on $v$ deliver exactly one pebble to $r$ along a minimising path.
Consequently every $r$-unsolvable distribution satisfies $D(v)<d_r(v)$ at
every vertex, and $\pi(G,r)$ is one more than the largest size of an
unsolvable distribution in that finite box.

\begin{lemma}\label{lem:factor}
$\pi(G,1)=\pi(G,2)=\pi(G,3)=3$, and hence $\pi(G)=3$.
\end{lemma}

\begin{proof}
For $r=1$ the box is $D(1)=0$, $D(2)\le1$, $D(3)\le2$, six distributions.
Exactly one of them is solvable, namely $D(2)=1$, $D(3)=2$: move $3\to2$ and
then $2\to1$. In the other five no winning play exists --- in
$(0,0,0),(0,1,0),(0,0,1),(0,1,1)$ no move is even legal, since every weight
is at least $2$, and from $(0,0,2)$ the only legal move is $3\to2$ at cost
$2$, which lands in $(0,1,0)$. The largest unsolvable distributions are
$(0,1,1)$ and $(0,0,2)$, of size $2$, so $\pi(G,1)=3$.

For $r=2$ the box is $D(1)\le1$, $D(2)=0$, $D(3)\le1$; no move is legal in
any of those four distributions, the largest being $(1,0,1)$, so
$\pi(G,2)=3$.

For $r=3$ the box is $D(1)\le2$, $D(2)\le1$, $D(3)=0$. Its largest member
$(2,1,0)$ is solvable: move $1\to2$ at cost $2$, reaching $(0,2,0)$, then
$2\to3$ at cost $2$. Of the distributions of size $2$, from $(2,0,0)$ the
only legal move is $1\to2$, which lands in $(0,1,0)$ and sticks, and $(1,1,0)$
admits no legal move; both are unsolvable, so $\pi(G,3)=3$. Now
\eqref{eq:unrooted} gives $\pi(G)=3$.
\end{proof}

\begin{lemma}\label{lem:dist}
In $G\times H$ with root $(x,y)$ one has
$d_{(x,y)}\big((i,j)\big)=d_x(i)\,d_y(j)$.
\end{lemma}

\begin{proof}
The weights along a path from $(i,j)$ to $(x,y)$ split into those of its
$G$-direction steps and those of its $H$-direction steps; the two projections
are walks from $i$ to $x$ in $G$ and from $j$ to $y$ in $H$, so the product
is at least $d_x(i)d_y(j)$. Concatenating a minimising path in the second
coordinate with one in the first attains that value.
\end{proof}

From now on $H=G$ and the root is $(1,1)$. By Lemma~\ref{lem:dist},
\[
\begin{array}{c|ccccccccc}
 v & (1,1) & (1,2) & (1,3) & (2,1) & (2,2) & (2,3) & (3,1) & (3,2) & (3,3)\\\hline
 d(v) & 1 & 2 & 3 & 2 & 4 & 6 & 3 & 6 & 9\\
 D(v) & 0 & 0 & 0 & 0 & 0 & 1 & 0 & 1 & 7
\end{array}
\]
where $D$ is the distribution we exhibit: seven pebbles on $(3,3)$ and one
each on $(2,3)$ and $(3,2)$, so $|D|=9$. The unweighted skeleton of
$G\times G$ is the $3\times3$ rook graph $K_3\,\square\,K_3$ on $9$ vertices
and $18$ edges; in the labelling $k\mapsto(\lfloor k/3\rfloor+1,\,k\bmod3+1)$,
$0\le k\le8$, its graph6 string is \texttt{H\{S\{aSf}.

\begin{theorem}\label{thm:main}
$D$ is not $(1,1)$-solvable. Consequently
$\pi(G\times G,(1,1))\ge10>9=\pi(G,1)\pi(G,1)$.
\end{theorem}

\begin{proof}
Put $\varphi(v)=1/d(v)$ and $\VAL(E)=\sum_vE(v)\varphi(v)$ for a
distribution $E$ on $G\times G$. Since $d(u)\le w(u,u')d(u')$ for every edge,
we have $\varphi(u')\le w(u,u')\varphi(u)$, so a move along $(u,u')$ changes
$\VAL$ by
\[
  -\bigl(w\varphi(u)-\varphi(u')\bigr)=:-\mathrm{loss}(u,u')\le0 .
\]
Thus $\VAL$ never increases. It does not by itself settle the matter, since
$\VAL(D)=7/9+1/6+1/6=10/9$ exceeds $1$; the obstruction is integrality, and
the rest of the proof is an accounting of it.

Suppose some play from $D$ reaches $(1,1)$, and stop it the instant a pebble
lands there. The stopping distribution $E$ has $E(1,1)=1$, so
\[
  \tfrac{10}{9}=\VAL(D)=\VAL(E)+L=1+R+L,
\]
where $L$ is the total loss of the play and
$R:=\sum_{v\neq(1,1)}E(v)\varphi(v)$; whence
\begin{equation}\label{eq:budget}
  R+L=\tfrac19 .
\end{equation}
Both terms are nonnegative, so every single move of the play has loss at most
$1/9$.

By Lemma~\ref{lem:dist} the loss of a move that changes the first coordinate
$i\to i'$ while the second is $j$ equals $\ell(i\to i')/d_1(j)$, where
$\ell(i\to i')=w(i,i')/d_1(i)-1/d_1(i')$ is computed in $G$ alone; and
symmetrically. The six values of $\ell$ are
\begin{gather*}
  \ell(2\to1)=\ell(3\to1)=0,\qquad \ell(3\to2)=\tfrac16,\qquad
  \ell(2\to3)=\tfrac23,\\
  \ell(1\to2)=\tfrac32,\qquad \ell(1\to3)=\tfrac83 .
\end{gather*}
Dividing by $d_1(j)\in\{1,2,3\}$, the moves of loss at most $1/9$ are exactly
these sixteen: the twelve moves in which a coordinate drops to $1$ (loss $0$);
$(2,3)\to(2,2)$ and $(3,2)\to(2,2)$, of loss $1/12$; and $(3,3)\to(2,3)$ and
$(3,3)\to(3,2)$, of loss $1/18$. Every other move costs at least
$1/6>1/9$, so none occurs in the play. Since off the root
$\varphi\ge\varphi(3,3)=1/9$, the quantity $R$ is either $0$ or at least
$1/9$, and \eqref{eq:budget} leaves two cases.

\emph{Case B: $L=0$ and $R=1/9$.} Only the twelve lossless moves occur, and
each of them lands on a vertex having a coordinate equal to $1$; in
particular none lands on $(3,2)$. So $(3,2)$ never gains a pebble, and its
single pebble can never be spent, every weight being at least $2$. Hence
$E(3,2)\ge1$ and $R\ge\varphi(3,2)=1/6>1/9$, a contradiction.

\emph{Case A: $R=0$ and $L=1/9$.} The positive losses available are $1/18$
and $1/12$ (the others exceeding the budget), and $1/9=4/36$ while
$1/18=2/36$ and $1/12=3/36$; the only multiset summing to $4/36$ is
$\{2/36,2/36\}$. So the play uses exactly two moves of loss $1/18$ and no
move of loss $1/12$, and every other move it makes is lossless; call the
fourteen moves of loss $0$ or $1/18$ \emph{available}. We now count pebbles at
four vertices; $R=0$ means every vertex other than $(1,1)$ is empty at the
end.

No available move lands on $(3,3)$: such a move raises a coordinate to $3$
and so costs at least $\ell(2\to3)/d_1(3)=2/9$. Hence the seven pebbles on
$(3,3)$ are only spent, in moves of weight $3$ (the two lossless moves
$(3,3)\to(1,3)$ and $(3,3)\to(3,1)$) and of weight $2$ (the two moves of loss
$1/18$, namely $(3,3)\to(2,3)$ and $(3,3)\to(3,2)$). As exactly two of the
latter occur, $3a+2\cdot2=7$, so $a=1$: precisely one weight-$3$ move leaves
$(3,3)$, delivering one pebble to $(1,3)$ or to $(3,1)$.

Say $p$ of the two lossy moves go to $(2,3)$ and $q=2-p$ to $(3,2)$. The only
available move landing on $(2,3)$ is $(3,3)\to(2,3)$, so $(2,3)$ receives
exactly $p$ pebbles and must spend all $1+p$ of them, in moves of weight $2$
or $3$. A single pebble can pay for no move, so $1+p\neq1$ and $p\ge1$;
symmetrically $q\ge1$; hence
$p=q=1$. So $(2,3)$ and $(3,2)$ each hold exactly two pebbles, which they can
only spend on their weight-$2$ moves $(2,3)\to(1,3)$ and $(3,2)\to(3,1)$.

Counting what arrives at $(1,3)$ and $(3,1)$: one pebble from the weight-$3$
move out of $(3,3)$, one at $(1,3)$ from $(2,3)$ and one at $(3,1)$ from
$(3,2)$, three in all, split $2+1$ or $1+2$ between the two vertices. But from
$(1,3)$, the only available move is the weight-$3$ move $(1,3)\to(1,1)$: the
alternative $(1,3)\to(1,2)$ has loss $\ell(3\to2)/d_1(1)=1/6$, and moves
raising the first coordinate cost at least $1/2$. The same holds at $(3,1)$.
With at most two pebbles on each, no move is possible, so those three pebbles
survive and $R\ge3\cdot\tfrac13=1>0$, a contradiction.

Both cases being impossible, no play from $D$ reaches $(1,1)$. As
$|D|=9$, this gives $\pi(G\times G,(1,1))\ge10$, while
$\pi(G,1)\pi(G,1)=9$ by Lemma~\ref{lem:factor}.
\end{proof}

\begin{corollary}\label{cor:three}
Conjectures~\ref{conj:W} and~\ref{conj:Wst}, in the form quoted above from
the e-print source, are false, and so is the unrooted statement
\eqref{eq:WG}; the answer to Sieben's question is no.
\end{corollary}

\begin{proof}
Theorem~\ref{thm:main} contradicts Conjecture~\ref{conj:W} at
$G=H$, $x=y=1$. Conjecture~\ref{conj:Wst} implies
Conjecture~\ref{conj:W} at $s=t=1$, so it fails too. For \eqref{eq:WG},
Lemma~\ref{lem:factor} gives $\pi(G)=\pi(H)=3$, whence
$\pi(G\times G)\ge\pi(G\times G,(1,1))\ge10>9=\pi(G)\pi(H)$. Note that the
last step needs $\pi(G,v)=3$ at \emph{every} $v$, not only at the root of the
witness: \eqref{eq:WG} is weaker than Conjecture~\ref{conj:W}, and a rooted
counterexample alone would not refute it.
\end{proof}

\section{The exact value}

A terminating exhaustive computation in integer arithmetic over the
$46656$-state box described above, re-run by the program that accompanies this
note, gives $\pi(G\times G,(1,1))=10$: no distribution of ten pebbles is
$(1,1)$-unsolvable. This equality, unlike the lower bound, is not proved by
hand, and neither Theorem~\ref{thm:main} nor Corollary~\ref{cor:three}
depends on it.

\section{Status}

\noindent\textbf{What is untouched.} The weights of $G$ are not all
equal, so nothing here bears on Graham's conjecture, on the unweighted rooted
conjecture $\pi(G\times H,(x,y))\le\pi(G,x)\pi(H,y)$, or on the unweighted
$st$-conjecture. Chung's proved region is not contradicted either: in a
weighted power of $K_2$ with parallel weights equal, every edge lies on a
minimum-product path to the root, that is, $d(u)=w(u,u')d(u')$ for the
endpoint $u$ further from the root. In our $G$ that fails at the edge
$\{2,3\}$, where $d_1(3)=3<4=w(2,3)d_1(2)$ and $d_1(2)=2<6=w(2,3)d_1(3)$, so
the edge is geodesic in neither direction and $G$ is not such a cube.

\medskip\noindent\textbf{What is not claimed.} No minimality: we do not
assert that $G\times G$ is a counterexample of least order, of least total
weight, or of least maximum weight, among weighted graphs or even among
weighted triangles. No priority either: we make no claim about what has or
has not appeared in print.

\medskip\noindent\textbf{The wording of the conjectures.} Both statements are
quoted from the e-print source of \cite{HHH}, where they are unnumbered; the
printed text of \emph{Discrete Math.} \textbf{312} (2012) 2286--2293 was not
accessible to us, so the printed statements and their numbering are
unverified here. What is refuted is accordingly the e-print form: we cannot
exclude that the printed versions carry a hypothesis on the weights that the
e-print statements do not. Our $G$ is tight in the sense
recalled in Section~1, so a tightness hypothesis on the factors would be met;
a hypothesis
requiring every edge to be geodesic to the root would not, as the paragraph on
Chung's region above records.

\section{Verification}

Everything a reader needs in order to redo the work is printed above: the
three weights of $G$, the definition of the product, the distribution $D$,
the distance vector, and a graph6 encoding of the skeleton. The hand proof of
Theorem~\ref{thm:main} uses no computer.

A short program accompanies this note, in Python and using only its standard
library, which reads those printed data and re-derives the quantities claimed
here: it rebuilds $G$ and $G\times G$ and checks the graph6 string against
them; recomputes $\pi(G,v)$ at all three roots and $\pi(G)$; confirms that
$D$ is unsolvable twice over, once by enumerating all $229$ distributions
reachable from $D$ with no use of distances, boxes or potentials, and once by
mechanising the proof above --- every move of a successful play costs at most
$1/9$, only sixteen of the thirty-six directed moves do, and the $112$ states
reachable from $D$ by those sixteen never carry a pebble on $(1,1)$;
recomputes $\pi(G\times G,(1,1))=10$; and re-derives every fraction of the
loss table and of the case analysis exactly, with no floating point. It also
runs controls in which the same semantics reproduce pebbling numbers
published by other authors. All $60$ of its checks passed, and every value
printed here agrees with the value it recomputed; the transcript records in
addition some checks on objects not discussed in this note. The program and a
transcript of that run accompany this note; neither is needed in order to redo
the hand proof, since every input it consumes is printed above.

\begin{thebibliography}{10}

\bibitem{Chung}
F.~R.~K. Chung,
\emph{Pebbling in hypercubes},
SIAM J. Discrete Math. \textbf{2} (1989), no.~4, 467--472.

\bibitem{Hurlbert}
G.~H. Hurlbert,
\emph{General graph pebbling},
Discrete Appl. Math. \textbf{161} (2013), no.~9, 1221--1231.
\href{https://doi.org/10.1016/j.dam.2012.03.010}{doi:10.1016/j.dam.2012.03.010}.

\bibitem{HHH}
D.~S. Herscovici, B.~D. Hester, and G.~H. Hurlbert,
\emph{Generalizations of Graham's pebbling conjecture},
Discrete Math. \textbf{312} (2012), no.~15, 2286--2293.
\href{https://doi.org/10.1016/j.disc.2012.03.032}{doi:10.1016/j.disc.2012.03.032}.
All quotations and line references are to the source of the e-print
\href{https://arxiv.org/abs/0905.3197}{arXiv:0905.3197}.

\bibitem{Sieben}
N.~Sieben,
\emph{A graph pebbling algorithm on weighted graphs},
J. Graph Algorithms Appl. \textbf{14} (2010), no.~2, 221--244.
\href{https://doi.org/10.7155/jgaa.00205}{doi:10.7155/jgaa.00205};
\href{https://arxiv.org/abs/0904.1651}{arXiv:0904.1651}.

\end{thebibliography}

\end{document}
